Consider the case of a vortex within the fluid. Strain within the fluid's
flow can naturally stretch it. Incompressibility forces the tube to narrow
as it lengthens, squeezing its fluid particles inward. For a locally circular
tube, the length \(\ell\) and radius \(r\) respond to a positive axial
strain \(\sigma\) as
\[
 \dot\ell=\sigma\ell>0,
 \qquad r\propto\ell^{-1/2}.
\]
As the radius shrinks, the circulation \(\Gamma\) enclosed by the tube
determines how fast the fluid turns around it. Its mean swirl speed \(v\) is
\[
 v=\frac{|\Gamma|}{2\pi r}.
\]
If circulation persists strongly enough through that narrowing, the fluid
circulates faster around a smaller core.

Viscosity diffuses momentum between neighbouring regions, reducing velocity
differences and spreading the rotation through the surrounding fluid.
It can carry rotation out of the narrowing tube and reduce its enclosed
circulation. The vorticity equation contains this competition between
stretching and diffusion:
\[
 \partial_t\omega+(u\cdot\nabla)\omega
 =(\omega\cdot\nabla)u+\nu\Delta\omega,
 \qquad \omega=\nabla\times u.
\]
While the tube's radius is shrinking, diffusion spreads rotation over
a distance
\[
 r_{\mathrm{diff}}\sim\sqrt{\nu t}.
\]
The material tube can keep narrowing as its fluid particles move inward,
even while rotation spreads beyond its boundary and its enclosed
circulation weakens.
